% F09：本书原创矢量示意；依据所在节的定义、证明或明确模型。
\begin{figure}[htbp]
\centering
\begin{tikzpicture}[x=1cm,y=1cm,>=stealth,every node/.style={font=\small},line width=.65pt]
\draw[->,gray] (-0.12,0)--(5.6,0) node[right] {\(x\)};\draw[->,gray] (0,-0.12)--(0,2.7);
\node[] at (2.6,3.4) {\(f=\mathbf1_{(-1,1)}\) 的中心平均};
\draw[AnalysisBlue,very thick](.7,0)--(.7,1.8)--(2.7,1.8)--(2.7,0);\fill(3.7,0)circle(2pt)node[below]{\(2\)};\draw[orange!80!black,thick](.7,-.5)--(6.7,-.5);\node[] at (3.6,-0.9) {半径 \(3\)，平均值 \(1/3\)};
\draw[->,gray] (9.38,0)--(12.5,0) node[right] {\(x\)};\draw[->,gray] (9.5,-0.12)--(9.5,2.8);
\node[] at (9.5,3.4) {中心极大函数};
\draw[AnalysisBlue,thick] plot coordinates {(6.9500,0.5750) (6.9746,0.5792) (6.9993,0.5835) (7.0239,0.5878) (7.0486,0.5922) (7.0732,0.5966) (7.0978,0.6011) (7.1225,0.6057) (7.1471,0.6104) (7.1717,0.6151) (7.1964,0.6199) (7.2210,0.6248) (7.2457,0.6298) (7.2703,0.6348) (7.2949,0.6399) (7.3196,0.6451) (7.3442,0.6504) (7.3688,0.6558) (7.3935,0.6613) (7.4181,0.6668) (7.4428,0.6725) (7.4674,0.6782) (7.4920,0.6841) (7.5167,0.6900) (7.5413,0.6961) (7.5659,0.7022) (7.5906,0.7085) (7.6152,0.7149) (7.6399,0.7214) (7.6645,0.7280) (7.6891,0.7347) (7.7138,0.7416) (7.7384,0.7486) (7.7630,0.7557) (7.7877,0.7630) (7.8123,0.7704) (7.8370,0.7779) (7.8616,0.7856) (7.8862,0.7935) (7.9109,0.8015) (7.9355,0.8097) (7.9601,0.8180) (7.9848,0.8266) (8.0094,0.8353) (8.0341,0.8441) (8.0587,0.8532) (8.0833,0.8625) (8.1080,0.8720) (8.1326,0.8817) (8.1572,0.8916) (8.1819,0.9017) (8.2065,0.9121) (8.2312,0.9227) (8.2558,0.9335) (8.2804,0.9446) (8.3051,0.9560) (8.3297,0.9677) (8.3543,0.9796) (8.3790,0.9919) (8.4036,1.0044) (8.4283,1.0173) (8.4529,1.0305) (8.4775,1.0441) (8.5022,1.0580) (8.5268,1.0723) (8.5514,1.0870) (8.5761,1.1021) (8.6007,1.1176) (8.6254,1.1336) (8.6500,1.1500)};
\draw[AnalysisBlue,thick] plot coordinates {(10.3500,1.1500) (10.3746,1.1336) (10.3993,1.1176) (10.4239,1.1021) (10.4486,1.0870) (10.4732,1.0723) (10.4978,1.0580) (10.5225,1.0441) (10.5471,1.0305) (10.5717,1.0173) (10.5964,1.0044) (10.6210,0.9919) (10.6457,0.9796) (10.6703,0.9677) (10.6949,0.9560) (10.7196,0.9446) (10.7442,0.9335) (10.7688,0.9227) (10.7935,0.9121) (10.8181,0.9017) (10.8428,0.8916) (10.8674,0.8817) (10.8920,0.8720) (10.9167,0.8625) (10.9413,0.8532) (10.9659,0.8441) (10.9906,0.8353) (11.0152,0.8266) (11.0399,0.8180) (11.0645,0.8097) (11.0891,0.8015) (11.1138,0.7935) (11.1384,0.7856) (11.1630,0.7779) (11.1877,0.7704) (11.2123,0.7630) (11.2370,0.7557) (11.2616,0.7486) (11.2862,0.7416) (11.3109,0.7347) (11.3355,0.7280) (11.3601,0.7214) (11.3848,0.7149) (11.4094,0.7085) (11.4341,0.7022) (11.4587,0.6961) (11.4833,0.6900) (11.5080,0.6841) (11.5326,0.6782) (11.5572,0.6725) (11.5819,0.6668) (11.6065,0.6612) (11.6312,0.6558) (11.6558,0.6504) (11.6804,0.6451) (11.7051,0.6399) (11.7297,0.6348) (11.7543,0.6298) (11.7790,0.6248) (11.8036,0.6199) (11.8283,0.6151) (11.8529,0.6104) (11.8775,0.6057) (11.9022,0.6011) (11.9268,0.5966) (11.9514,0.5922) (11.9761,0.5878) (12.0007,0.5835) (12.0254,0.5792) (12.0500,0.5750)};
\draw[AnalysisBlue,thick] plot coordinates {(8.6500,2.3000) (10.3500,2.3000)};
\draw[AnalysisBlue,fill=white](8.65,2.3)circle(1.7pt);\fill[AnalysisBlue](8.65,1.15)circle(1.7pt);\draw[AnalysisBlue,fill=white](10.35,2.3)circle(1.7pt);\fill[AnalysisBlue](10.35,1.15)circle(1.7pt);
\end{tikzpicture}
\caption[移动平均与极大函数的超水平集]{移动平均与极大函数的超水平集。对中心极大算子，\(Mf\) 在开区间内为一，在 \(|x|\ge1\) 时为 \(1/(1+|x|)\)。端点值是二分之一；非中心版本不能套用这条曲线。}
\label{b1:fig:visual-f09}
\end{figure}
